% Anexos de la tesis.
% Incluye tablas, instrumentos, evidencias u otros materiales de soporte.

\section*{ANEXOS}
\addcontentsline{toc}{section}{ANEXOS}

\subsection*{Anexo A: Transcripción de entrevista con operador}
\addcontentsline{toc}{subsection}{Anexo A: Transcripción de entrevista con operador}

\textit{Metadatos:} fecha 12 de abril de 2026; entrevistado \textbf{Raúl Laura Anncasi} (operador de flota); entrevistador Roger Infa Sánchez (investigador); empresa \textbf{Empresa de Transportes Characato -- Sabandia C-8 S.A.}; duración aproximada 18 minutos.

\vspace{0.5em}
\noindent\textbf{Investigador:} Gracias por su tiempo. Para el estudio necesito entender cómo trabaja usted día a día y qué información le falta para coordinar las unidades.

\vspace{0.35em}
\noindent\textbf{Raúl Laura Anncasi:} Con gusto. Yo soy quien asigna qué bus sale con qué chofer y en qué ruta, según el turno. El problema es que muchas veces el usuario en paradero no sabe si el bus ya viene o si se quedó atrás, y nosotros tampoco tenemos una pantalla única donde ver en vivo dónde está cada unidad. Eso justifica que el sistema permita \textbf{ver en el mapa los buses activos por ruta}: si yo puedo filtrar por línea y ver los puntos, reduzco llamadas y malentendidos (relacionado con HU-001).

\vspace{0.35em}
\noindent\textbf{Investigador:} ¿Qué pasa cuando un pasajero no conoce bien el paradero?

\vspace{0.35em}
\noindent\textbf{Raúl Laura Anncasi:} Pasa mucho. La gente pregunta ``¿cuál es el más cercano?'' porque en la calle hay varios puntos. Si la aplicación puede usar el GPS del celular y decir cuál es el \textbf{paradero más cercano}, eso baja la incertidumbre (relacionado con HU-002).

\vspace{0.35em}
\noindent\textbf{Investigador:} ¿Y si el usuario quiere planificar antes de salir de casa?

\vspace{0.35em}
\noindent\textbf{Raúl Laura Anncasi:} Claro. Muchos quieren ver la \textbf{ruta completa}, el trazo y el orden de paraderos, no solo el nombre. Por eso tiene sentido poder \textbf{consultar una ruta con su trazo y paraderos ordenados} (relacionado con HU-003).

\vspace{0.35em}
\noindent\textbf{Investigador:} ¿Qué información esperan en el paradero?

\vspace{0.35em}
\noindent\textbf{Raúl Laura Anncasi:} Lo mínimo es saber cuánto falta para el próximo bus. Hoy eso es muy subjetivo: depende de radio o de la experiencia del cobrador. Si el sistema puede mostrar un \textbf{ETA} en el paradero seleccionado, el usuario deja de adivinar (relacionado con HU-004).

\vspace{0.35em}
\noindent\textbf{Investigador:} Desde operación, ¿qué necesitan ustedes registrar?

\vspace{0.35em}
\noindent\textbf{Raúl Laura Anncasi:} Necesito registrar \textbf{asignación bus--conductor--ruta por turno}. Si no queda registrado, después no sabemos quién manejaba qué unidad si hay un incidente (relacionado con HU-005). También necesitamos que solo los buses en \textbf{servicio activo} estén publicando ubicación; si está en mantenimiento, no debe aparecer como disponible en el mapa (relacionado con HU-007).

\vspace{0.35em}
\noindent\textbf{Investigador:} ¿Cómo llega la ubicación al sistema?

\vspace{0.35em}
\noindent\textbf{Raúl Laura Anncasi:} Tiene que haber un flujo para \textbf{actualizar la ubicación GPS en tiempo real} cuando el bus está en ruta. Si eso no es estable, todo lo demás falla: mapa, ETA, todo (relacionado con HU-006).

\vspace{0.35em}
\noindent\textbf{Investigador:} ¿Y del lado administrativo?

\vspace{0.35em}
\noindent\textbf{Raúl Laura Anncasi:} Las rutas cambian de color en plano, se corrigen trazos, se agregan paraderos. Necesitamos \textbf{registrar y editar rutas} con su geometría, y \textbf{paraderos} con ubicación bien puesta. Igual con \textbf{buses y conductores}: datos base tienen que estar limpios o el monitoreo se ensucia (relacionado con HU-008, HU-009 y HU-010).

\vspace{0.35em}
\noindent\textbf{Investigador:} ¿Hay algo más que consideren para mejorar tiempos?

\vspace{0.35em}
\noindent\textbf{Raúl Laura Anncasi:} Sí. Si más adelante podemos mirar \textbf{historial de tráfico por tramo y franja}, eso ayuda a explicar por qué un día el bus tarda más; también sirve para comparar antes y después en la investigación (relacionado con HU-011 y HU-012).

\vspace{0.35em}
\noindent\textbf{Investigador:} ¿Algo que deba quedar explícito como criterio de calidad?

\vspace{0.35em}
\noindent\textbf{Raúl Laura Anncasi:} Que lo que se muestre al usuario sea \textbf{consistente} con lo que operamos. Si el mapa dice que el bus está en un punto, tiene que ser verificable; si no, perdemos confianza. Por eso las historias deben venir con criterios de aceptación claros.

\vspace{0.35em}
\noindent\textbf{Investigador:} ¿Autoriza usar esta entrevista como insumo para el levantamiento de requerimientos?

\vspace{0.35em}
\noindent\textbf{Raúl Laura Anncasi:} Sí, autorizo, siempre que no se publiquen datos personales innecesarios de los choferes más allá de lo operativo.

\subsection*{Anexo B: Cuestionario de Diagnóstico y Aceptación de la Propuesta (Pre-test)}
\addcontentsline{toc}{subsection}{Anexo B: Cuestionario de Diagnóstico y Aceptación de la Propuesta (Pre-test)}

Este instrumento fue diseñado y administrado de forma virtual (Google Forms) como diagnóstico inicial para validar la problemática de los tiempos de espera y evaluar la viabilidad de adopción de la plataforma tecnológica TransiGo (Formulario en línea disponible en: \url{https://docs.google.com/forms/d/e/1FAIpQLSd4nm3DCIx5Zsz994Qbhek1Fln9DLuQV4g2ISjdefbInDQCag/viewform?usp=sharing&ouid=110680708735571704289}).

\begin{enumerate}
    \item \textbf{Sección I: Datos Generales}
    \begin{itemize}
        \item ¿Cuál es su rango de edad? (15--24 años, 25--34 años, 35--49 años, 50 años a más).
        \item ¿Cuál es su género? (Masculino, Femenino, Prefiero no decirlo).
        \item ¿Cuál es su ocupación principal? (Estudiante, Trabajador dependiente, Trabajador independiente, Ama/o de casa, Otro).
        \item ¿Con qué frecuencia utiliza el transporte público en Arequipa? (Todos los días, 3 a 5 veces por semana, 1 a 2 veces por semana, Ocasionalmente).
    \end{itemize}
    
    \item \textbf{Sección II: Problemática Actual del Transporte}
    \begin{itemize}
        \item En promedio, ¿cuánto tiempo espera en el paradero hasta que llegue su bus? (Menos de 5 minutos, Entre 5 y 15 minutos, Entre 15 y 30 minutos, Más de 30 minutos).
        \item ¿Con qué frecuencia no sabe cuándo llegará el bus que espera? (Siempre, Casi siempre, A veces, Casi nunca).
        \item ¿Ha perdido compromisos importantes debido a la impuntualidad del transporte público? (Sí, frecuentemente; Sí, alguna vez; No).
        \item ¿Qué problema le genera mayor incomodidad al usar el transporte público? (No saber cuándo llegará el bus, Tiempos de espera muy largos, Rutas poco claras o desconocidas, Saturación en el paradero, Otro).
    \end{itemize}

    \item \textbf{Sección III: Uso de Tecnología}
    \begin{itemize}
        \item ¿Dispone de un teléfono con acceso a internet? (Sí, No).
        \item ¿Utiliza actualmente alguna aplicación para consultar rutas o movilizarse? (Sí, con frecuencia; Sí, pero pocas veces; No, no conozco ninguna; No, prefiero no usarlas).
        \item ¿Cómo calificaría su nivel de facilidad para usar aplicaciones en el celular? (Muy fácil, Regular, Difícil, No uso aplicaciones).
    \end{itemize}

    \item \textbf{Sección IV: Aceptación de la Solución Propuesta}
    \begin{itemize}
        \item Si existiera una aplicación que mostrara en tiempo real la ubicación de los buses y el tiempo estimado de arribo (ETA), ¿la utilizaría? (Definitivamente sí, Probablemente sí, No estoy seguro/a, Probablemente no).
        \item ¿Qué información le resultaría más útil? (ETA, Ubicación del bus en tiempo real, Paradero más cercano, Ruta completa del bus, Aviso de retrasos/desvíos).
        \item ¿A través de qué dispositivo preferiría acceder a esta información? (Celular - navegador web, App descargable, Pantallas informativas, Indistinto).
        \item ¿Estaría dispuesto/a a compartir su ubicación GPS para recibir información personalizada? (Sí, sin problema; Sí, solo si es necesario; No, me preocupa mi privacidad).
        \item ¿Qué tan importante considera que Arequipa cuente con un sistema de transporte en tiempo real? (Escala 1 al 5: Nada importante a Muy importante).
    \end{itemize}
    
    \item \textbf{Sección V: Comentarios Adicionales}
    \begin{itemize}
        \item (Opcional) ¿Tiene alguna sugerencia sobre cómo debería funcionar esta aplicación? (Pregunta abierta).
    \end{itemize}
\end{enumerate}

\subsection*{Anexo C: Cuestionario de Evaluación Post-Implementación (Post-test)}
\addcontentsline{toc}{subsection}{Anexo C: Cuestionario de Evaluación Post-Implementación (Post-test)}

Cuestionario administrado de forma virtual a los 20 usuarios que participaron en la simulación del experimento, con el propósito de evaluar la usabilidad, la precisión percibida en el entorno simulado del sistema TransiGo y validar cualitativamente la reducción de los tiempos de espera simulados (Formulario en línea disponible en: \url{https://docs.google.com/forms/d/e/1FAIpQLScsQj0BwQLlRjwv0YloabRAw33XlVJ2R0XSc1pA5M8pu55-sA/viewform}).

\begin{enumerate}
    \item \textbf{Sección I: Datos de Identificación}
    \begin{itemize}
        \item ID de Usuario (Asociado al registro de la simulación del experimento).
    \end{itemize}
    
    \item \textbf{Sección II: Evaluación de la Plataforma TransiGo}
    \begin{itemize}
        \item ¿Considera que el uso de la aplicación TransiGo redujo efectivamente su tiempo de espera en el paradero durante la simulación de rutas? (Sí, significativamente; Sí, levemente; No, el tiempo de espera fue el mismo).
        \item En promedio, ¿cuántos minutos estima que ahorró por viaje utilizando la aplicación en el escenario simulado? (Menos de 5 minutos, Entre 5 y 10 minutos, Entre 10 y 20 minutos, Más de 20 minutos).
        \item ¿Con qué frecuencia logró acudir al paradero de manera sincronizada con la llegada del bus simulada en la plataforma? (Siempre, Casi siempre, A veces, Casi nunca).
        \item ¿Cómo califica la precisión del tiempo estimado de llegada (ETA) que mostró la aplicación en la simulación? (Excelente, Buena, Regular, Mala).
        \item ¿El sistema web redujo su nivel de estrés o incertidumbre respecto a los tiempos de arribo en la prueba simulada? (Sí, significativamente; Sí, levemente; No).
        \item ¿Qué funcionalidad de TransiGo le resultó de mayor utilidad práctica en el simulador? (Visualización del bus en tiempo real en el mapa, Tiempo estimado de llegada - ETA, Localización del paradero más cercano, Trazo de rutas).
        \item ¿Cómo califica la facilidad de uso e interactividad de la interfaz del panel web simulado? (Muy fácil, Fácil, Regular, Difícil).
        \item Califique de manera general su nivel de satisfacción con el sistema implementado en la simulación: (Escala del 1 al 5: Muy insatisfecho a Muy satisfecho).
        \item ¿Recomendaría la implementación a gran escala de esta plataforma en el Sistema Integrado de Transportes (SIT) de Arequipa? (Definitivamente sí, Probablemente sí, No estoy seguro/a, Definitivamente no).
        \item (Opcional) Escriba sugerencias o comentarios para futuras versiones de la aplicación TransiGo: (Pregunta abierta).
    \end{itemize}
\end{enumerate}

